\chapter{测试自动化与硬件在环}
\label{chap:test-automation-hil}

嵌入式测试需要同时验证纯软件逻辑、内核接口、完整镜像和真实电气行为。单一测试环境无法覆盖所有故障，因此应建立从主机单元测试到硬件在环的分层体系，并让每一层具有明确的发现目标。

\section{测试金字塔}

推荐层次包括：
\begin{enumerate}
  \item 主机单元测试：算法、解析器、状态机和错误分类；
  \item 仿真与模拟：QEMU、FVP、虚拟外设、fake ioctl 和软件故障注入；
  \item 内核测试：KUnit、kselftest、LTP 和驱动回归；
  \item 板级接口测试：真实 GPIO、总线、DMA、存储和网络；
  \item HIL 系统测试：电源、复位、外部信号、温度和长时间运行；
  \item 生产测试：覆盖装配、器件和身份写入，不重复全部研发测试。
\end{enumerate}

越靠近硬件，测试速度越慢、成本越高且故障定位越困难。应把可脱离硬件的逻辑尽量下沉到快速测试，但不能用 mock 通过替代目标板电气验证。

\section{SIL、模型在环与虚拟平台}

模型在环（Model-in-the-Loop, MIL）通常让控制算法模型与被控对象或环境模型在同一建模环境中闭环运行，用于在生成或编写产品代码前验证算法。软件在环（Software-in-the-Loop, SIL）则把生成或手写的软件实现编译为主机进程，并由环境模型或测试替身驱动。QEMU 和 FVP 进一步提供目标 ISA、异常、中断与虚拟设备环境\cite{qemu-system-doc,arm-fvp-overview}。这些名词在不同行业的边界略有差异，项目应直接说明被测对象、执行 ISA、时间模型和被替代的硬件，而不是只写缩写。

推荐按以下顺序上移同一组场景：

\begin{enumerate}
  \item 用主机单元测试验证纯计算和状态不变量；
  \item 用虚拟时钟及设备替身验证超时、乱序、重试和资源耗尽；
  \item 在 QEMU/FVP 中验证目标 ABI、启动、异常和完整镜像；
  \item 在目标板上验证真实控制器、DMA、缓存和性能；
  \item 在 HIL 中闭合电气输入、物理对象和故障反应链路。
\end{enumerate}

每上移一层都应复用测试向量、事件 ID 和判定不变量。若每个环境拥有完全独立的脚本与判断逻辑，测试之间无法形成证据链，环境差异也容易被误认为产品差异。

虚拟平台用例应保存模型版本、非默认参数、镜像摘要、DTB/ACPI、随机种子和完整串口。涉及超时时，要区分被测系统的虚拟时间与 CI 主机的墙上时间：前者验证产品语义，后者负责发现模型卡死或基础设施失联。

\section{契约测试与差分测试}

硬件抽象层、驱动与服务之间的接口应具有可在多个环境执行的契约测试。例如一个传感器读取接口不只测试正常值，还应固定验证：短传输、CRC 错误、超时、设备复位、时间戳回退、重复样本和取消请求。

差分测试让同一输入在参考实现、优化实现、模拟器和目标板上运行，再比较稳定语义。适合比较的对象包括：

\begin{itemize}
  \item 协议编码结果、状态转移和错误分类；
  \item 启动阶段、固件 ABI 和设备枚举不变量；
  \item DSP/控制算法在已定义定点误差范围内的输出；
  \item 文件系统或数据库操作后的业务不变量；
  \item 异常原因、恢复动作和安全输出。
\end{itemize}

不应直接比较线程交错、无约束日志顺序、虚拟平台周期数或浮点最后一位。测试 oracle 必须允许规范定义的自由度，同时拒绝会破坏产品契约的差异。

\section{可测试架构}

驱动访问、时间、随机数、文件系统和网络应通过窄接口注入，使状态机能在主机上控制输入。测试替身必须模拟错误、短读写、超时和乱序，而不只是返回成功。

协议解析器应采用长度受限的纯函数，适合单元测试和 fuzzing。生产状态机应能使用虚拟时钟推进，避免测试依赖真实长延时。

\section{内核测试}

KUnit 适合测试内核内部纯逻辑和边界条件；kselftest 验证用户可见内核 ABI；LTP 覆盖广泛系统调用和内核功能。BSP 不需要盲目运行全部测试，但应选择与配置和产品功能相关的集合，并记录跳过原因。

驱动测试应覆盖 probe defer、资源获取失败、IRQ 超时、runtime PM、解绑和模块反复加载。使用 fault injection、kmemleak、KASAN、KCSAN 和 lockdep 时，要区分诊断配置引入的时序变化。

\section{HIL 架构}

典型 HIL 由控制主机、可编程电源、继电器或电子开关、串口采集、信号发生器、测量仪器和被测设备组成：

\begin{center}
\begin{tikzpicture}[node distance=6mm]
  \tikzset{hil/.style={draw=bookline,fill=bookpaper,rounded corners=1mm,
    minimum width=32mm,minimum height=10mm,align=center,font=\small}}
  \node[hil] (host) {测试控制器\\调度与记录};
  \node[hil,right=of host] (fixture) {电源/继电器\\信号注入};
  \node[hil,right=of fixture] (dut) {DUT};
  \node[hil,below=of fixture] (measure) {串口、示波器\\电流与网络};
  \draw[->,bookteal] (host)--(fixture);
  \draw[<->,bookteal] (fixture)--(dut);
  \draw[<->,bookteal] (host)--(measure);
  \draw[<->,bookteal] (measure)--(dut);
\end{tikzpicture}
\end{center}

控制协议应返回机器可读状态和时间戳。测试程序不能只等待固定时间再搜索日志，应等待可观测条件并设置总超时。

\section{测试清单与资产关联}

每次测试必须关联：软件镜像摘要、源码版本、配置、硬件版本、DUT 序列、工装版本、仪器校准状态、环境条件和原始日志。只有测试名称和 Pass/Fail 不能支持复现。

\begin{lstlisting}[language=C,caption={测试结果 manifest 的逻辑字段}]
struct test_run_record {
    char run_id[40];
    char image_digest[65];
    char hardware_revision[16];
    char fixture_revision[16];
    char environment_id[32];
    uint64_t started_at;
    uint32_t passed;
    uint32_t failed;
    uint32_t infrastructure_errors;
};
\end{lstlisting}

基础设施错误必须与产品失败区分，例如串口采集断开、继电器未动作、仪器超量程和网络存储不可用。否则不稳定工装会污染产品质量数据。

\section{电源循环与掉电注入}

HIL 应能直接切断 DUT 主电源，并测量实际下降曲线。软件 reboot、PMIC orderly shutdown 和拔除外部适配器可能具有完全不同的掉电行为。

随机掉电测试需要记录切断时对应的业务阶段，并在恢复后执行文件系统和业务不变量检查。测试分布应覆盖擦除、写入、fsync、rename、数据库提交和 OTA 切换等关键窗口。

\section{信号与总线故障注入}

可通过模拟开关、可编程负载或协议代理注入断线、短路、抖动、延迟、CRC 错误和设备无响应。注入设备必须保证默认状态安全，并在控制主机崩溃时释放危险输出。

总线分析仪看到的线上事务应与驱动 trace 和应用事务 ID 关联。这样可以区分硬件未响应、驱动未发起和应用超时策略错误。

\section{长稳与统计}

长稳测试应持续记录内存、句柄、温度、功耗、错误计数、重连次数和最大延迟。只在结束时检查设备是否仍在线，会遗漏逐步退化。

性能结果应保存分布、最大值和测试条件。若系统要求极低失效率，需要根据样本量与运行时间说明统计置信度，不能以少量“零失败”直接证明目标可靠性。

\section{工装校准与自检}

HIL 本身也是测量系统。电压、电流、时间和温度通道需要校准；继电器和插座有寿命；线缆与探头会改变信号。每次批量测试前应使用参考负载或 loopback 执行工装自检。

工装固件、脚本和接线图必须版本化。修改工装后需要确认历史结果是否仍可比较。

\section{CI 与实验室调度}

硬件资源有限时，CI 应把快速主机测试作为每次提交门禁，把目标板 smoke test 用于合并请求，把长稳、温度和破坏性故障注入安排为周期任务。调度器需要设备锁、超时回收、自动重新上电和隔离故障板。

未经确认不应自动把任意开发分支烧写到共享生产工装。镜像来源、签名和允许的设备池应受到访问控制。

虚拟平台任务还应限制 CPU、内存、磁盘和许可证并发数。模型启动失败、许可证不可用和 worker 资源耗尽属于基础设施错误，不应记为 DUT 功能失败。重复执行只适合确认基础设施瞬态问题，不能用“重跑直到通过”掩盖竞态或时序缺陷。

\section{验收清单}

\begin{itemize}
  \item 每种测试层是否具有明确目标，避免重复和空白？
  \item 同一场景能否跨主机、虚拟平台、目标板和 HIL 复用输入与不变量？
  \item mock 是否覆盖错误、短传输、超时和并发，而不仅是成功路径？
  \item 虚拟时间、墙上时间、随机种子和模型配置是否可追溯？
  \item HIL 能否区分 DUT 失败与基础设施失败？
  \item 测试结果是否关联镜像、硬件、工装、环境和原始日志？
  \item 掉电与信号注入是否达到真实硬件故障边界？
  \item 工装是否具有校准、自检、默认安全状态和寿命管理？
  \item 长稳结果是否保存趋势、最大值和统计条件？
\end{itemize}
